\documentclass[11pt]{article}
\usepackage[margin=1in]{geometry}
\usepackage[T1]{fontenc}
\usepackage[utf8]{inputenc}
\usepackage{lmodern}
\usepackage{microtype}
\usepackage{amsmath,amssymb,amsthm,mathtools}
\usepackage{enumitem}
\usepackage{booktabs}
\usepackage{array}
\usepackage{longtable}
\usepackage{hyperref}
\usepackage{xcolor}
\hypersetup{colorlinks=true,linkcolor=blue,citecolor=blue,urlcolor=blue}
\setlist[itemize]{leftmargin=1.5em}
\setlist[enumerate]{leftmargin=1.75em}

\title{A Conceptual Primer on Learning Beyond Pure Chaining\\\large Credal Beliefs, Representation Growth, and Cognitive Synergy\\\large Dated synthesis: 2026-04-17}
\author{}
\date{}

\newtheorem{definition}{Definition}
\newtheorem{principle}{Principle}
\newtheorem{claim}{Claim}
\newtheorem{remark}{Remark}
\newcommand{\wm}{world-model}
\newcommand{\wmpln}{WM-PLN}
\newcommand{\pln}{PLN}
\newcommand{\R}{\mathbb{R}}
\newcommand{\Bel}{\mathrm{Bel}}
\newcommand{\VOI}{\mathrm{VOI}}

\begin{document}
\maketitle

\begin{abstract}
This note develops a conceptual answer to a central question in symbolic
probabilistic reasoning: what should a system do when pure chaining no longer
recovers the right belief state?  The answer proposed here is not to abandon
reasoning, but to embed it in a richer learning architecture.  The key move is
to recognize that additive or local chaining succeeds only on a restricted
fragment, while realistic partially observable environments force one to track
beliefs over hidden structure, ambiguity over models, and the value of
constructing better internal representations.

The note argues for a five-part picture.  First, pure chaining must be given a
sharp semantic boundary.  Second, hidden-state ambiguity should be represented
credally rather than collapsed prematurely to a single scalar.  Third, learning
should search for better sufficient statistics, better latent variables, and
better compression maps rather than merely tuning local rule weights.  Fourth,
limited computation should be managed by attention and value-of-information
policies.  Fifth, these mechanisms should compose in a cognitively synergistic
way: symbolic chaining, Bayesian belief tracking, imprecise belief
representation, representation learning, and attention control should each
repair the others' blind spots instead of competing for a single universal
formalism.
\end{abstract}

\tableofcontents

\section{Purpose and scope}

This article is a companion to a broader primer on world-model Markov, hidden
Markov, and POMDP reasoning.  Its focus is narrower and more strategic:
\begin{quote}
If pure symbolic probabilistic chaining does not fully solve the reasoning
problem, where should learning enter?
\end{quote}

The goal is not to replace semantic reasoning with statistical heuristics.
Rather, the aim is to describe a layered architecture in which learning enters
precisely where chaining reaches principled limits.

The note is written as a conceptual primer for future work and future
conversations.  It therefore emphasizes design commitments, failure modes, and
research directions rather than implementation detail.

\section{The boundary of pure chaining}

\subsection{What pure chaining is good at}

Pure chaining is strongest in settings where:
\begin{itemize}
  \item the semantic model is explicit,
  \item local composition rules are exact under known side conditions,
  \item or the posterior state admits an additive sufficient statistic.
\end{itemize}

In such cases, local symbolic or evidence-carrying reasoning can be exact,
transparent, and computationally attractive.

\subsection{What breaks}

The boundary appears when the system is asked to reason through hidden
variables, partial observability, or underdetermined latent structure.  Then
one encounters at least four kinds of failure.
\begin{enumerate}
  \item \textbf{Non-identifiability:} different latent models induce the same
        observed law.
  \item \textbf{State uncertainty:} what matters for action is a posterior over
        hidden states, not a directly observed symbol frequency.
  \item \textbf{Representation mismatch:} the currently chosen predicates or
        symbols are not sufficient statistics for the real task.
  \item \textbf{Resource limits:} even when the semantics is clear, exact
        inference may be too expensive to run exhaustively.
\end{enumerate}

\subsection{The methodological lesson}

\begin{principle}[Boundary honesty]
When pure chaining ceases to be semantically faithful, the right response is not
to stretch its formulas beyond their domain.  The right response is to change
the representation, the uncertainty object, the computational policy, or all
three.
\end{principle}

This point is subtle.  It does not say symbolic reasoning is obsolete.  It says
symbolic reasoning must be embedded in a richer architecture that can learn
better representations, maintain ambiguity honestly, and choose what to compute
next.

\section{Belief state rather than derived edge}

\subsection{The first conceptual shift}

The first shift beyond pure chaining is from \emph{derived edge confidence} to
\emph{belief state}.  In a partially observable system, the operative question
is usually not:
\begin{quote}
What scalar truth value should I attach to this next symbolic consequence?
\end{quote}
but rather:
\begin{quote}
What posterior state over hidden structure should guide action, prediction, and
further inference?
\end{quote}

This is the fundamental contribution of filtering in hidden-state models and of
belief-state semantics in POMDPs.

\subsection{The second shift}

The second shift is from \emph{single posterior state} to \emph{set of possible
posterior states}.  When observations underdetermine latent structure, belief
itself is ambiguous.  A single posterior number can be a lossy and misleading
summary.

\section{Credal beliefs as the first learning-aware response}

\subsection{Why credal semantics comes first}

The most direct response to non-identifiability is a credal one.  If multiple
hidden models remain compatible with the same observations, then the hidden
belief should be represented as a set or interval, not as a point estimate
chosen too early.

\begin{definition}[Credal latent belief]
Given a family of compatible latent models $\Theta$ and an observation history
$h$, the credal belief over a latent query $q$ is the interval
\[
\left[\inf_{\theta \in \Theta} P_\theta(q \mid h),\;
\sup_{\theta \in \Theta} P_\theta(q \mid h)\right].
\]
\end{definition}

\subsection{Why this matters for symbolic systems}

This move is important because it turns a no-go result into a usable semantic
object.

Positive example:
\begin{itemize}
  \item if two hidden-state models agree on all observations but disagree
        maximally on a latent-state query, the correct answer is the full
        interval from $0$ to $1$.
\end{itemize}

Negative example:
\begin{itemize}
  \item a single scalar ``confidence'' may hide the fact that the posterior is
        structurally ambiguous rather than merely noisy.
\end{itemize}

\subsection{A conceptual interpretation of confidence}

In this setting, confidence is no longer best interpreted as a primitive extra
field attached to every inference step.  A deeper interpretation is:
\begin{quote}
Confidence reflects the width, stability, or fragility of the compatible belief
set induced by the current observations and assumptions.
\end{quote}

This connects credal semantics to the older symbolic intuition that truth value
should include both strength and confidence, but it grounds that distinction in
hidden-model ambiguity rather than in an ad hoc numeric convention.

\section{Learning as sufficient-statistic discovery}

\subsection{The basic idea}

When the current observable symbols are not sufficient for the task, one should
learn derived observables, summary statistics, or latent predicates that make
the problem more identifiable or more decision-relevant.

\begin{principle}[Representation before weight fitting]
When pure chaining fails because the currently visible vocabulary is
insufficient, the first learning move should be to improve the representation,
not merely to retune local rule weights.
\end{principle}

\subsection{What this means concretely}

This can take several forms.
\begin{itemize}
  \item Learn history features that better predict latent mode or future value.
  \item Learn abstract predicates that summarize long observational traces.
  \item Learn context-sensitive sufficient statistics for special subdomains.
  \item Learn regime detectors that decide which local reasoning laws are safe
        in the present situation.
\end{itemize}

\subsection{Why this is better than ``fixing the formulas''}

If the failure comes from omitted structure, then stretching a local chaining
formula will not recover what the representation never recorded.  The proper
repair is to make the missing structure visible, either explicitly or through a
better learned surrogate.

\section{Learning as latent-variable construction}

\subsection{From hidden-state recovery to latent-state invention}

Another conceptual step is to stop assuming that the true latent state space is
already known.  In many real problems, the best internal state is not recovered
from the world but constructed by the agent to improve prediction and action.

\begin{definition}[Constructed latent state]
A constructed latent state is an internal variable introduced because it
compresses history in a way that improves future prediction, control, or
explanation, even if it does not correspond one-to-one with a unique external
hidden variable.
\end{definition}

\subsection{Why this matters}

This move changes the problem from:
\begin{quote}
Which latent state is really there?
\end{quote}
to:
\begin{quote}
Which internal latent variables make reasoning and action more effective?
\end{quote}

That reframing is especially important under persistent non-identifiability.
One need not solve metaphysical uniqueness in order to construct a useful
belief-state architecture.

\subsection{A key consequence}

If latent variables are allowed to be constructed rather than merely recovered,
then learning enters not only as parameter estimation but also as ontology
growth.  The system may decide that its current symbolic vocabulary is too poor
for the task and extend it.

\section{Learning as attention and belief compression}

\subsection{The computational problem}

Even when the correct belief state is well-defined, tracking it exactly may be
too expensive.  So another place for learning is not the semantics itself but
the \emph{compression policy} deciding which parts of belief to maintain
faithfully.

\subsection{Attention as a first-class epistemic policy}

An attention mechanism can be understood as a learned projection
\[
\Bel \mapsto \widehat{\Bel}
\]
that preserves action-relevant distinctions while forgetting the rest.

\begin{principle}[Value-sensitive compression]
A belief-compression map should be evaluated not by whether it preserves all
semantic distinctions, but by whether it preserves the distinctions that matter
for high-value decisions.
\end{principle}

\subsection{Why this is not merely approximation}

This is not just ``cheap inference.''  It is a principled architectural layer:
the system learns what latent distinctions are worth maintaining under limited
time and memory.  In that sense, attention is a learned epistemic policy over
belief state.

\section{Learning as meta-inference}

\subsection{Beliefs about filters}

One can go a step further and treat the update rule itself as uncertain.  Then
learning no longer tunes only model parameters; it compares alternative
inference procedures, representation choices, and latent constructions.

\subsection{Why this is attractive}

This matters because real systems are often uncertain not just about the world
but about:
\begin{itemize}
  \item which model family is appropriate,
  \item which latent vocabulary is worth using,
  \item which summary statistics are sufficient,
  \item and which approximation policy best trades off cost and fidelity.
\end{itemize}

\subsection{The conceptual payoff}

Learning can then be read as model selection and update-rule selection at the
same time.  A reasoning system is not merely updating beliefs inside a fixed
formal shell; it is learning which shell best supports competent behavior.

\section{A cognitively synergistic architecture}

\subsection{Why no single mechanism is enough}

Pure chaining fails on hidden structure.  Exact Bayesian filtering is often too
expensive or too brittle under model mismatch.  Credal semantics preserves
ambiguity but does not on its own invent better latent states.  Representation
learning may discover useful features but can lose interpretability.  Attention
compression manages cost but risks forgetting the wrong thing.

The right picture is therefore not replacement but composition.

\subsection{The synergy stack}

A coherent learning-aware world-model architecture might include:
\begin{enumerate}
  \item \textbf{Exact symbolic/evidence reasoning} on fragments where the model
        really supports it.
  \item \textbf{Belief-state tracking} for hidden-state semantics.
  \item \textbf{Credal belief representation} when observations leave latent
        structure underdetermined.
  \item \textbf{Representation learning} to discover more useful summaries,
        predicates, or latent variables.
  \item \textbf{Attention and compression} to allocate finite computation toward
        action-relevant distinctions.
  \item \textbf{Meta-level selection} over models, filters, and reasoning modes.
\end{enumerate}

\subsection{The conceptual moral}

\begin{principle}[Cognitive synergy]
The role of learning is not merely to replace symbolic reasoning with a more
flexible black box.  Its deeper role is to supply missing representations,
missing statistics, and missing control policies exactly where symbolic chaining
reaches principled semantic limits.
\end{principle}

\section{Where learning should enter in world-model reasoning}

The following decomposition is especially useful.

\subsection{Before chaining}

Learning can enter before chaining by improving the representation:
\begin{itemize}
  \item feature discovery,
  \item latent construction,
  \item regime detection,
  \item and vocabulary growth.
\end{itemize}

\subsection{During chaining}

Learning can enter during chaining by controlling:
\begin{itemize}
  \item which rules to attempt,
  \item which contexts to reveal,
  \item which belief dimensions to maintain,
  \item and when to abstain.
\end{itemize}

\subsection{After chaining}

Learning can enter after chaining by revising:
\begin{itemize}
  \item which model family remains credible,
  \item which latent variables are worth preserving,
  \item and which approximation policy is acceptable under observed error.
\end{itemize}

\section{What this implies for POMDP-capable reasoning}

\subsection{The planning connection}

In partially observable control, the natural hidden-state object is the belief
state.  A symbolic world-model system that wants to reason honestly in this
setting should therefore:
\begin{itemize}
  \item track or approximate filtering and smoothing objects,
  \item carry ambiguity credally when hidden-state models are underdetermined,
  \item and learn belief compressions or latent refinements that improve action
        quality.
\end{itemize}

\subsection{The practical implication}

Learning beyond pure chaining is therefore not optional decoration.  It is the
mechanism that turns a symbolic reasoning substrate into an architecture capable
of operating under partial observability, ambiguity, and finite resources.

\subsection{The conservative universal-prediction leg}

One useful new consequence is that the POMDP-facing story is no longer only
about posterior semantics.  Along any fixed action stream, a controlled hidden
Markov law can now be compared against a concrete controlled Solomonoff-style
semimeasure obtained by lifting the ordinary observation-alphabet universal
mixture and ignoring actions.  This is deliberately conservative, but it is
already enough to connect ambiguity-aware belief tracking to finite-horizon
log-loss guarantees.

The conceptual lesson is important:
\begin{itemize}
  \item universal prediction can be attached first at the level of fixed action
        streams;
  \item this already yields meaningful regret guarantees for partially
        observable control;
  \item and the remaining gap is no longer ``whether a bridge exists'' but
        whether one wants a genuinely action-aware universal mixture rather
        than the current conservative comparator.
\end{itemize}

The lower-semicomputability hypothesis in that theorem is also conceptually
honest.  A raw hidden-state parameter object with arbitrary probability
measures is not automatically computable.  So a learning architecture that
wants to reason with universal predictors should expose a computable or
lower-semicomputable model subclass explicitly instead of suppressing the
domain condition.

\section{A theorem-shaped research program}

The following research directions are especially coherent.

\begin{enumerate}
  \item \textbf{Credal latent belief theorems.}  Formalize lower and upper
        filtering and smoothing beliefs for families of hidden models.
  \item \textbf{Computable hidden-model classes.}  Package model families whose
        induced observation laws are lower semicomputable, so universal
        prediction bounds apply without extra ad hoc hypotheses.
  \item \textbf{Representation-growth semantics.}  State what it means for a
        learned latent variable or predicate to improve identifiability or value.
  \item \textbf{Belief-compression theorems.}  Prove conditions under which a
        learned compressed belief state preserves action ranking or bounded loss.
  \item \textbf{Value-of-information policies.}  Formalize when revealing more
        context is worth its cost.
  \item \textbf{Mode-switching semantics.}  Specify when the system should rely
        on exact chaining, posterior inference, credal intervals, or learned
        compression.
  \item \textbf{Meta-filter learning.}  Formalize the selection of update rules
        or model families based on predictive performance and decision quality.
\end{enumerate}

\section{A compact primer for future conversations}

Future work can take the following as a stable conceptual baseline.

\begin{enumerate}
  \item Pure chaining is exact only on a restricted semantic fragment.
  \item Hidden-state reasoning requires belief states, not merely derived local
        supports.
  \item Non-identifiability should be represented credally, not hidden.
  \item Learning should improve representations before merely retuning rule
        parameters.
  \item Constructed latent variables are often more important than recovery of a
        uniquely ``true'' hidden state.
  \item Attention is best understood as a learned value-sensitive belief
        compression policy.
  \item The most powerful architecture is synergistic: symbolic reasoning,
        Bayesian belief tracking, credal semantics, representation learning,
        attention, and meta-level control each repair the others' failure modes.
  \item In POMDP-like settings, learning beyond pure chaining is the route from
        local symbolic inference to competent action under ambiguity.
  \item A conservative controlled Solomonoff comparator already gives a first
        universal-prediction anchor for that route, even before a fully
        action-aware universal mixture exists.
\end{enumerate}

\section{Conclusion}

The right response to the limitations of pure chaining is not surrender, and it
is not unlimited confidence in larger learned models.  The right response is a
clean architecture in which symbolic reasoning, probabilistic belief tracking,
imprecise uncertainty, learned representation growth, and attention control are
all given distinct and principled roles.

From this perspective, learning enters exactly where it should:
\begin{itemize}
  \item to discover better sufficient statistics,
  \item to build more useful latent variables,
  \item to preserve ambiguity honestly,
  \item to compress belief under resource limits,
  \item and to choose which reasoning mode should be active at all.
\end{itemize}

That picture preserves the strongest virtues of symbolic world-model reasoning
while acknowledging the realities of hidden structure, partial observability,
and finite computation.  It therefore provides a coherent foundation for
learning-aware \wmpln{} and POMDP-capable world-model reasoning.

\end{document}
